\section{Assiomi dello Schema in Turtle}
\label{sec:app_turtle}

Lo schema \`e memorizzato in RDF/XML (\texttt{videogames.owl}). Gli estratti seguenti
lo riportano in Turtle, pi\`u leggibile, per mostrare come la sintassi DL delle sezioni
precedenti diventa un insieme di triple RDF. Restrizioni, intersezioni e liste sono
rappresentate da nodi blank (\texttt{[ ... ]} e \texttt{( ... )}). Si omettono
\texttt{rdfs:label} e \texttt{rdfs:comment}.

\subsection*{Classe definita con restrizione esistenziale}

\begin{lstlisting}[style=turtle]
vg:PlayStationGame a owl:Class ;
    rdfs:subClassOf vg:VideoGame ;
    owl:equivalentClass [ a owl:Class ;
        owl:intersectionOf ( vg:VideoGame
            [ a owl:Restriction ;
              owl:onProperty vg:availableOn ;
              owl:someValuesFrom vg:PlayStationPlatform ] ) ] .
\end{lstlisting}

$\texttt{PlayStationGame} \equiv \texttt{VideoGame} \sqcap \exists\,\texttt{availableOn.PlayStationPlatform}$.
L'intersezione \`e una lista RDF di due elementi: la classe \texttt{VideoGame} e
una restrizione anonima.

\subsection*{Classe definita con datatype restriction}

\begin{lstlisting}[style=turtle]
vg:HighlyRatedGame a owl:Class ;
    rdfs:subClassOf vg:VideoGame ;
    owl:equivalentClass [ a owl:Class ;
        owl:intersectionOf ( vg:VideoGame
            [ a owl:Restriction ;
              owl:onProperty vg:metacriticScore ;
              owl:someValuesFrom [ a rdfs:Datatype ;
                  owl:onDatatype xsd:integer ;
                  owl:withRestrictions ( [ xsd:minInclusive 85 ] ) ] ] ) ] .
\end{lstlisting}

Il filler di \texttt{someValuesFrom} non \`e una classe ma un \emph{data range}:
il sottoinsieme di \texttt{xsd:integer} vincolato dal facet \texttt{xsd:minInclusive 85}.

\subsection*{Classe definita con complemento}

\begin{lstlisting}[style=turtle]
vg:StandaloneGame a owl:Class ;
    rdfs:subClassOf vg:VideoGame ;
    owl:equivalentClass [ a owl:Class ;
        owl:intersectionOf ( vg:VideoGame
            [ a owl:Class ;
              owl:complementOf [ a owl:Restriction ;
                  owl:onProperty vg:sequelOf ;
                  owl:someValuesFrom vg:VideoGame ] ] ) ] .
\end{lstlisting}

$\texttt{StandaloneGame} \equiv \texttt{VideoGame} \sqcap \lnot\exists\,\texttt{sequelOf.VideoGame}$.

\subsection*{Condizione solo necessaria con cardinalit\`a}

\begin{lstlisting}[style=turtle]
vg:GameWithSingleDeveloper a owl:Class ;
    rdfs:subClassOf vg:VideoGame ,
        [ a owl:Restriction ;
          owl:onProperty vg:developedBy ;
          owl:maxCardinality "1"^^xsd:nonNegativeInteger ] .
\end{lstlisting}

Due assiomi \texttt{rdfs:subClassOf} distinti, nessun \texttt{owl:equivalentClass}:
\[ \texttt{GameWithSingleDeveloper} \sqsubseteq \texttt{VideoGame} \sqcap {\le}1\,\texttt{developedBy} \]
L'assenza di \texttt{owl:onClass} rende la cardinalit\`a non qualificata.

\subsection*{Propriet\`a transitiva e inversa}

\begin{lstlisting}[style=turtle]
vg:sequelOf a owl:ObjectProperty , owl:TransitiveProperty ;
    rdfs:domain vg:VideoGame ;
    rdfs:range vg:VideoGame ;
    owl:inverseOf vg:hasSequel .

vg:hasSequel a owl:ObjectProperty ;
    rdfs:domain vg:VideoGame ;
    rdfs:range vg:VideoGame .
\end{lstlisting}

\subsection*{Property chain}

\begin{lstlisting}[style=turtle]
vg:sharesDeveloperWith a owl:ObjectProperty , owl:SymmetricProperty ;
    rdfs:domain vg:VideoGame ;
    rdfs:range vg:VideoGame ;
    owl:propertyChainAxiom ( vg:developedBy vg:developerOf ) .
\end{lstlisting}

La catena \`e una lista ordinata: l'ordine delle propriet\`a conta
($\texttt{developedBy} \circ \texttt{developerOf}$).

\subsection*{Disjointness e chiavi}

\begin{lstlisting}[style=turtle]
[] a owl:AllDisjointClasses ;
    owl:members ( vg:ConsolePlatform vg:PCPlatform vg:MobilePlatform ) .

vg:Review owl:hasKey ( vg:reviewedGame vg:reviewerName ) .
\end{lstlisting}

\texttt{owl:AllDisjointClasses} dichiara disgiunte a coppie tutte le classi della lista
con un solo assioma, invece di $n(n-1)/2$ assiomi \texttt{owl:disjointWith}.
